\documentclass[pageno]{jpaper}

\newcommand{\IWreport}{Fall 2019}
\newcommand{\quotes}[1]{``#1''}


\widowpenalty=9999

\usepackage[normalem]{ulem}

\begin{document}

\title{Accelerating Entropy-Based Transformer Calibration}

\author{Michael Hu \\Adviser: Prof. Karthik Narasimhan}

\date{}
\maketitle

\thispagestyle{empty}
\doublespacing
\begin{abstract}
This document is intended to serve as a sample you can use for independent work reports.  We provide some guidelines on content and formatting.  They are not required, but they might be helpful.
\end{abstract}

\section{Introduction}

This document is a sample you can use for formatting your independent work
report.  An example template with all bib and supporting files can be downloaded
from the IW website at {\em http://iw.cs.princeton.edu}.

\section{Preparation Instructions}

\subsection{Paper Formatting}
\label{section:formatting}

There are no minimum or maximum length limits on IW reports.  
We are including this template because we think it will be helpful
for citing things properly and for including figures into formatted
text.  If you are using \LaTeX~\cite{lamport94} 
to typeset your paper, then we strongly suggest
that you start from the template available at
http://iw.cs.princeton.edu -- this
document was prepared with that template.  
If you are using a different software package to typeset your paper, 
then you can still use this document as a reasonable sample of 
how your report might look.  Table~\ref{table:formatting} is a suggestion
of some formatting guidelines, as well as being an example of how to
include a table in a Latex document.

\begin{table}[hbt]
  \centering
  \begin{tabular}{|l|l|}
    \hline
    \textbf{Field} & \textbf{Value}\\
    \hline
    \hline
    Paper size & US Letter 8.5in $\times$ 11in\\
    \hline
    Top margin & 1in\\
    \hline
    Bottom margin & 1in\\
    \hline
    Left margin & 1in\\
    \hline
    Right margin & 1in\\
    \hline
    Body font & 12pt\\
    \hline
    Abstract font & 12pt, italicized\\
    \hline
    Section heading font & 14pt, bold\\
    \hline
    Subsection heading font & 12pt, bold\\
    \hline
  \end{tabular}
  \caption{Formatting guidelines. }
  \label{table:formatting}
\end{table}

\textbf{Please ensure that you include page numbers with your
submission}. This makes it easier for readers to refer to
different parts of your paper when they provide comments.

We highly recommend you use bibtex for managing your references and citations.  You can add bib entries to a references.bib file throughout the semester (e.g., as you read papers) and then they will be ready for you to cite when you start writing the report.  If you use bibtex, please note that the references.bib file provided in the template example includes some format-specific incantations at the top of the file.  If you substitute your own bib file, you will probably want to include these 
incantations at the top of it.

\subsection{Citations and Footnotes}

There are various reasons to cite prior work and include it as references in your bibliography.  For example, If you are improving upon 
prior work, you should include
a full citation for the work in the bibliography \cite{nicepaper,nicepaper2}. 
You can also cite information that is used as background or
explanation\cite{Salzberg:2005}.  In addition to citing scholarly papers or books, you can
also create bibtex entries for webpages or other sources.  Many online
databases allow you to download a premade bibtex entry for each paper
you access.  You can simply copy-paste these into your references.bib
file.

Sometimes you want to footnote something, such as a web
site.\footnote{http://www.cs.princeton.edu}  Note that the footnote
number comes after the punctuation.

\subsection{Figures and Tables.}

Figure \ref{fig:gray} shows an example of how to include a figure in
your report.  
Ensure that the figures and
tables are legible.  Please also ensure that you refer to your
figures in the main text. Make sure that your figures will be legible
in the expected forms that the report will be read.  If you expect someone
to print it out in gray-scale, then make sure the figures are legible 
when printed that way.  

\begin{figure}[hbt]
\centering
\includegraphics[width=0.75\linewidth]{gray.jpg}
\caption{This is a gray image.}
\label{fig:gray}
\end{figure}

In Section~\ref{section:formatting}, an example of a table was given.
(Note that the ``S'' in Section is capitalized.  Here's one more
example - see Table~\ref{table:data}.

\begin{table}[hbt]
  \centering
  \begin{tabular}{|l|l|} \hline
    \textbf{Some field} & \textbf{Another field}\\\hline
    200          &  10000 \\ \hline 
    400          &  20000 \\ \hline 
    800          &  40000 \\ \hline 
    1600        &  80000 \\ \hline 
    3200        &  160000 \\ \hline 
    6400        &  320000 \\ \hline 
  \end{tabular}
  \caption{Some data in a table. }
  \label{table:data}
\end{table}


Here's an example that shows how you can have side-by-side figures -
see Figure~\ref{fig:side-a} and Figure~\ref{fig:side-b}.  (Note that
the the ``F'' in Figure is capitalized. 

\begin{figure}[htbb]
\begin{minipage}[b]{0.5\linewidth}
\centering
\includegraphics[width=.75\linewidth]{checkerboard-squares-black-white.jpg}
\caption{Plain checkerboard.}
\label{fig:side-a}
\end{minipage}
\hspace{0.5cm}
\begin{minipage}[b]{0.5\linewidth}
\centering
\includegraphics[width=.75\linewidth]{swirl-squares-black-white.jpg}
\caption{Cool checkerboard.}
\label{fig:side-b}
\end{minipage}
\end{figure}

\subsection{Double Quotes.}

Latex double quotes are not the same as the double quote key on your
keyboard. The standard way of writing quotes and double quotes in
LaTeX is with `` and '' not with " and ".   

Now that may be confusing, so you may want to use the \textbackslash\{quotes\} command.  For
example \quotes{The quick brown fox.}



\subsection{Main Body.}

Avoid bad page or column breaks in
your main text, i.e., last line of a paragraph at the top of a
column or first line of a paragraph at the end of a column. If you
begin a new section or sub-section near the end of a column,
ensure that you have at least two (2)  lines of body text on the same
column. 

\section{Outline}  
The following is a possible outline for your paper.
\subsection{Introduction}
\begin{itemize}
\item Motivation and Goal (The goal of this project is...)
\item Overview of challenge and previous work 
\item Approach 
\item Summary of implementation
\item Summary of results
\end{itemize}

\subsection{Problem Background and Related Work}
\begin{itemize}
\item Survey of prior work with similar goals 
\item For each previous approach, explain what has been done and why it does not meet your goal
\end{itemize}

\subsection{Approach}
\begin{itemize}
\item Key novel idea
\item Why it is a good idea
\end{itemize}

\subsection{Implementation}
\begin{itemize}
\item System overview (flow chart of key steps?)
\item Subsection for each step or issue you addressed
\begin{itemize}
\item Problem statement
\item Possible approaches
\item Chosen approach and why
\item Implementaton details
\end{itemize}
\end{itemize}

\subsection{Evaluation}
\begin{itemize}
\item Experiment design...
\item Data...
\item Metrics...
\item Comparisons...
\item Qualitative results...
\item Quantitative results...
\end{itemize}

\subsection{Summary}
\begin{itemize}
\item Conclusions...
\item Limitations...
\item Future work...
\end{itemize}


\section{Ethics}

Your independent work report should abide by the basic standards of scholarly ethics and by the Princeton Honor Code. If you have any doubts about how to cite
other work, how to quote or include text or images from other works, or other issues, please discuss them with your project adviser or with the IW coordinators. 



\bstctlcite{bstctl:etal, bstctl:nodash, bstctl:simpurl}
\bibliographystyle{IEEEtranS}
\bibliography{references}

\end{document}

